\section{Chronological local factors}
\label{sec:peps_chronological_network}

Fix a chronological composition $W:\mathcal H_a\to\mathcal H_b$ as in
Definition~\ref{def:peps_source_circuit}. Let $G(W)$ be its finite set of gate
occurrences. At occurrence $g$, write $P_g\subseteq P$ for the actual
participants, $\Xi_g$ for its branch labels, and $c_g(\xi)$ for its original
coefficient. Equal gates at different times give different occurrences.
For a selection $A\subseteq P$, put
$T_A(W)=\{g:P_g\cap A\ne\varnothing\}$.
We use the restriction of a source-free composition to a set of owners and
the canonical grouping isometries of
Theorem~\ref{thm:peps_party_restriction} and
Definition~\ref{def:peps_party_partition}.
These are the local-factor identities used in
\cite[proof of Theorem~5.2]{OpenAI2026PolynomialPEPS}.

% This chapter does not assert the final equality between the sampled physical
% operator and a contraction of the party tensors. That identity is separate.

\begin{definition}[Full labels and their restrictions]
    \label{def:peps_chrono_full_labels}
    \lean{TNLean.PEPS.PairEffect.SourceCircuit.choiceAt_projectAllChoices,
        TNLean.PEPS.PairEffect.SourceCircuit.projectAllChoices,
        TNLean.PEPS.PairEffect.SourceCircuit.FullLabels,
        TNLean.PEPS.PairEffect.SourceCircuit.choiceAt_choicesOfFullLabels,
        TNLean.PEPS.PairEffect.SourceCircuit.choicesOfFullLabels,
        TNLean.PEPS.PairEffect.SourceCircuit.choicesOfFullLabels_eq_of_agree,
        TNLean.PEPS.PairEffect.SourceCircuit.choicesOfFullLabels_party_eq_of_agree,
        TNLean.PEPS.PairEffect.SourceCircuit.partySelector}
    \leanok
    \uses{def:peps_partial_source_choices}
    A full label is $\lambda\in\Lambda(W):=\prod_{g\in G(W)}\Xi_g$.
    Restricting it to the touched occurrences gives a partial choice
    $\pi_A\lambda\in\mathcal C_A(W)$. At every $g\in T_A(W)$, its recovered
    label is exactly $\lambda_g$. Thus $\pi_A\lambda=\pi_A\mu$ whenever
    $\lambda_g=\mu_g$ at all touched occurrences. In particular, the choice at
    one party $p$ depends only on the labels with $p\in P_g$.
    There is also a restriction $\pi_A^P:\mathcal C_P(W)\to\mathcal C_A(W)$
    from the actual all-party partial choices; it preserves the recovered label
    at every touched occurrence. These maps do not select labels at untouched
    occurrences.
\end{definition}

\begin{proof}
    \leanok
    Define the maps by following the chronology. Composition splits its occurrence
    set into the two original occurrence sets. At a gate, retain its given label
    if it is touched and the unique empty choice otherwise. Induction verifies
    the recovered labels and the asserted dependence.
\end{proof}

\begin{definition}[Labels at precisely the touched occurrences]
    \label{def:peps_chrono_touched_labels}
    \lean{TNLean.PEPS.PairEffect.SourceCircuit.allChoicesNonemptyLabelsEquiv,
        TNLean.PEPS.PairEffect.SourceCircuit.allChoicesNonemptyLabelsEquiv_apply,
        TNLean.PEPS.PairEffect.SourceCircuit.isTouched_all_iff,
        TNLean.PEPS.PairEffect.SourceCircuit.nonemptyGateLocations,
        TNLean.PEPS.PairEffect.SourceCircuit.nonemptyGateLocationsFintype,
        TNLean.PEPS.PairEffect.SourceCircuit.TouchedLabels,
        TNLean.PEPS.PairEffect.SourceCircuit.choiceAt_choicesOfTouchedLabels,
        TNLean.PEPS.PairEffect.SourceCircuit.choicesOfTouchedLabels,
        TNLean.PEPS.PairEffect.SourceCircuit.choicesOfTouchedLabels_choiceAt,
        TNLean.PEPS.PairEffect.SourceCircuit.choicesTouchedLabelsEquiv}
    \leanok
    \uses{def:peps_chrono_full_labels}
    The actual partial choices have the canonical description
    \begin{align}
        \mathcal C_A(W)\simeq\prod_{g\in T_A(W)}\Xi_g.
        \label{eq:peps_chrono_touched_labels}
    \end{align}
    The forward map recovers each original occurrence label, and the inverse
    inserts exactly these labels in the chronology. Write
    $G_+(W)=\{g:P_g\ne\varnothing\}$. This set is finite, and for the all-party
    selection, (\ref{eq:peps_chrono_touched_labels}) becomes
    $\mathcal C_P(W)\simeq\prod_{g\in G_+(W)}\Xi_g$.
    Empty-participant gates request no label, including when their branch set
    is empty.
\end{definition}

\begin{proof}
    \leanok
    The two maps are inverse by induction on the chronology. A gate is touched
    by the all-party selection exactly when it has a participant. Restricting
    the recovered labels to these occurrences gives the second equivalence.
\end{proof}

\begin{theorem}[The coefficient of an actual partial choice]
    \label{thm:peps_chrono_coefficient}
    \lean{TNLean.PEPS.PairEffect.SourceCircuit.coefficient_all_eq_prod_nonemptyGateLocations,
        TNLean.PEPS.PairEffect.SourceCircuit.coefficient_eq_prod_dite}
    \leanok
    \uses{def:peps_chrono_touched_labels}
    For every $\xi\in\mathcal C_A(W)$, with recovered labels $\xi_g$,
    \begin{align}
        c_A(\xi) & =\prod_{g\in T_A(W)}c_g(\xi_g), &
        c_P(\xi) & =\prod_{g\in G_+(W)}c_g(\xi_g).
        \label{eq:peps_chrono_coefficient}
    \end{align}
    Each gate coefficient occurs once. The aggregate scalars of empty-participant
    gates belong to the actual residual, rather than to this product.
\end{theorem}

\begin{proof}
    \leanok
    Induct on the chronology. Composition multiplies the two coefficient products;
    a touched gate contributes its original coefficient and an untouched gate
    contributes $1$. For the second identity, retain precisely the nonempty
    participant occurrences in the finite product.
\end{proof}

\begin{theorem}[Compatibility of owner restrictions]
    \label{thm:peps_chrono_restriction}
    \lean{TNLean.PEPS.PairEffect.Word.restrict_restrict_eval_heq_of_implies,
        TNLean.PEPS.PairEffect.Word.restrict_restrict_heq_of_implies,
        TNLean.PEPS.PairEffect.Word.restrict_withoutParty_atParty_eval_heq,
        TNLean.PEPS.PairEffect.Word.restrict_castInput_eval_heq,
        TNLean.PEPS.PairEffect.Word.restrict_comp_eval_heq,
        TNLean.PEPS.PairEffect.Word.restrict_frameList_eval_heq,
        TNLean.PEPS.PairEffect.Word.restrict_frame_eval_heq,
        TNLean.PEPS.PairEffect.Word.restrict_localPart_comp_eval_heq,
        TNLean.PEPS.PairEffect.Word.restrict_localPart_frame_eval_heq,
        TNLean.PEPS.PairEffect.Word.restrict_moveHead_eval_heq,
        TNLean.PEPS.PairEffect.Word.restrict_castLayouts_eval_heq}
    \leanok
    \uses{thm:peps_party_restriction}
    Let $U$ be a source-free composition. If $B\subseteq A$, successive restriction
    first to $A$ and then to $B$ gives the same composition as direct restriction
    to $B$, under the canonical identifications of the register lists. In
    particular, removing a party other than $p$ does not change the operator at
    $p$. Restriction and evaluation respect these identifications.

    Equality of restricted evaluations is preserved by composition, by adjoining
    the same spectator registers, and by the permutations that move one register
    past a fixed list. The same assertion holds for the source-free residuals of
    compositions with pair preparations, provided that the corresponding source
    register lists agree. The residual of a composition retains the later source
    registers as spectators while evaluating the earlier residual.
\end{theorem}

\begin{proof}
    \leanok
    Induct on the composition. A local map is kept precisely when its owner is
    selected. For a register exchange, inspect the membership of its two owners.
    This proves the equality of the restricted compositions themselves. The
    remaining assertions follow by composition of the restricted evaluations
    and by tensoring with the identity on the retained spectator memories.
    Identifying equal register lists leaves these operations unchanged.
\end{proof}

\begin{theorem}[Operations away from the selected parties]
    \label{thm:peps_chrono_disjoint}
    \lean{TNLean.PEPS.PairEffect.SourceCircuit.partialWord_all_gate_eq_of_touched,
        TNLean.PEPS.PairEffect.SourceCircuit.partialWord_all_gate_eq_of_untouched,
        TNLean.PEPS.PairEffect.SourceCircuit.restrict_partialWord_all_gate_localPart_eval_heq_id_of_disjoint,
        TNLean.PEPS.PairEffect.PreparedSourceGate.restrict_remaining_appendTail_eq_of_not_participating,
        TNLean.PEPS.PairEffect.PreparedSourceGate.restrict_remaining_appendTail_eval_eq_id_of_not_participating,
        TNLean.PEPS.PairEffect.Word.restrict_mapOwner_appendTail_eval_eq_id_of_disjoint,
        TNLean.PEPS.PairEffect.Word.restrict_localPart_mapOwner_appendTail_mapOwner_eval_eq_id_of_disjoint,
        TNLean.PEPS.PairEffect.Word.parties_mapOwner_of_sources_nil,
        TNLean.PEPS.PairEffect.Word.restrict_mapOwner_eval_eq_id_of_disjoint}
    \leanok
    \uses{thm:peps_chrono_restriction,def:peps_prepared_source_gate}
    Relabelling the owners of a source-free composition sends its set of
    participating parties to their image. If every relabelled owner lies outside
    $A$, its restriction to $A$ is the identity on the empty memory.
    After adjoining any spectator list, the restriction is the identity on the
    selected spectator memory. The same statement holds after a further
    relabelling and after moving all pair preparations to the beginning: source
    vectors may be arbitrary.

    Consequently, the residual of a prepared gate not incident to $p$ acts as
    the identity on the registers of $p$, independently of its branch label.
    At a nonempty-participant gate, the all-party partial choice selects exactly
    the corresponding prepared branch. At an empty-participant gate it retains
    the original aggregate operator, independently of the unique partial choice.
\end{theorem}

\begin{proof}
    \leanok
    Inspect the local maps, exchanges, and preparations in the composition.
    Every operation away from the selected owners restricts to the identity;
    the canonical permutations preserve the selected spectator order. The
    same induction after moving sources to the beginning proves the statement
    for the actual residual. The gate assertions then follow from the definition
    of the partial expansion and its touched-occurrence test.
\end{proof}

\begin{definition}[Complete chronological branches and their free sources]
    \label{def:peps_chrono_complete_branch}
    \lean{TNLean.PEPS.PairEffect.SourceCircuit.fullBranchWord,
        TNLean.PEPS.PairEffect.SourceCircuit.fullSourceVector,
        TNLean.PEPS.PairEffect.SourceCircuit.fixedSourceLayout,
        TNLean.PEPS.PairEffect.SourceCircuit.fullRemainingWord}
    \leanok
    \uses{def:peps_chrono_full_labels,def:peps_canonical_source_residual}
    For $\lambda\in\Lambda(W)$, replace every prepared gate by its indicated
    branch and write $W_\lambda$ for the resulting composition. Let $\mathcal S_W$
    be the ordered register list of all its pair preparations, with later
    preparations first, and let $s_\lambda\in\mathcal H_{\mathcal S_W}$ be their
    ordered product vector. The list $\mathcal S_W$ is independent of $\lambda$.
    There is an actual source-free composition
    $V_\lambda:\mathcal H_{\mathcal S_W\mathbin{+}a}\to\mathcal H_b$ obtained
    by moving these source registers to the beginning. In a composition of two
    chronological pieces, retain the later source registers while applying the
    earlier residual. A register retained as a spectator is moved past the
    free-source list before the residual acts.
\end{definition}

\begin{theorem}[Exact preparation of a complete branch]
    \label{thm:peps_chrono_complete_preparation}
    \lean{TNLean.PEPS.PairEffect.SourceCircuit.fullRemainingWord_gate_prepare_tmul,
        TNLean.PEPS.PairEffect.SourceCircuit.isAllowed_fullRemainingWord,
        TNLean.PEPS.PairEffect.SourceCircuit.sources_fullRemainingWord,
        TNLean.PEPS.PairEffect.SourceCircuit.eval_fullBranchWord_eq_fullRemainingWord,
        TNLean.PEPS.PairEffect.SourceCircuit.fullRemainingWord_gate_prepare,
        TNLean.PEPS.PairEffect.SourceCircuit.fullRemainingWord_prepare}
    \leanok
    \uses{def:peps_chrono_complete_branch,thm:peps_chrono_spectators}
    The composition $V_\lambda$ contains no pair preparation. If $W$ is allowed,
    then $V_\lambda$ is allowed. For every input vector $x$,
    \begin{align}
        W_\lambda x=V_\lambda(s_\lambda\otimes x),
        \label{eq:peps_chrono_complete_preparation}
    \end{align}
    where concatenated register memories are identified by the canonical
    associating isometry. This is an operator identity, and at a single gate
    it holds with arbitrary spectator input.
\end{theorem}

\begin{proof}
    \leanok
    For one gate use its exact prepared-branch identity and the compatibility of
    owner relabelling with concatenation. Elementary tensors span the gate and
    spectator memory. Induction on the chronology then proves
    (\ref{eq:peps_chrono_complete_preparation}); the composition step uses the
    later-first order of the source registers and the spectator identity. The
    source-free and allowedness assertions follow from the same construction.
\end{proof}

\begin{theorem}[Dependence of complete residuals on incident labels]
    \label{thm:peps_chrono_complete_locality}
    \lean{TNLean.PEPS.PairEffect.SourceCircuit.restrict_fullRemainingWord_eq_of_agree,
        TNLean.PEPS.PairEffect.SourceCircuit.restrict_fullRemainingWord_eq_of_choices_eq}
    \leanok
    \uses{thm:peps_chrono_disjoint,thm:peps_chrono_restriction,def:peps_chrono_full_labels}
    If two full labels $\lambda,\mu$ agree at every occurrence incident to $p$,
    then $V_\lambda|_p=V_\mu|_p$. Equivalently, equality of
    $\pi_{\{p\}}\lambda$ and $\pi_{\{p\}}\mu$ implies this equality of actual
    restricted operators.
\end{theorem}

\begin{proof}
    \leanok
    Induct on the chronology. At an incident gate the labels agree. At any other
    gate, Theorem~\ref{thm:peps_chrono_disjoint} identifies both restrictions with
    the spectator identity. Composition and retained registers respect these
    equalities by Theorem~\ref{thm:peps_chrono_restriction}. Recovering individual
    labels from the partial choice gives the equivalent formulation.
\end{proof}

\begin{definition}[Complete local residual maps]
    \label{def:peps_chrono_complete_local_maps}
    \lean{TNLean.PEPS.PairEffect.SourceCircuit.localRemainingMap,
        TNLean.PEPS.PairEffect.SourceCircuit.norm_localRemainingMap_le_one,
        TNLean.PEPS.PairEffect.SourceCircuit.restrict_fullRemainingWord_eq_localRemainingMap}
    \leanok
    \uses{thm:peps_chrono_complete_locality,thm:peps_chrono_complete_preparation}
    For a party $p$ and an incident choice $\eta\in\mathcal C_{\{p\}}(W)$,
    define $L_p(\eta)$ by restricting $V_\lambda$ for any full label extending
    $\eta$. If no such extension exists, set $L_p(\eta)=0$.
    Theorem~\ref{thm:peps_chrono_complete_locality} makes the extending case
    independent of the extension. For every genuine full label,
    $V_\lambda|_p=L_p(\pi_{\{p\}}\lambda)$.
    If $W$ is allowed, every $L_p(\eta)$ has norm at most $1$, including the
    zero value at an unextendable choice.
\end{definition}

\begin{proof}
    \leanok
    Use the incident-label dependence to compare any two extending full labels.
    For an extending choice, restriction preserves allowedness, hence gives a
    contraction; the other case is the zero map.
\end{proof}

\begin{theorem}[Canonical factors are the actual singleton restrictions]
    \label{thm:peps_chrono_canonical_factors}
    \lean{TNLean.PEPS.PairEffect.Word.groupByPartyIso_eval_eq_tensorRestrictions}
    \leanok
    \uses{thm:peps_party_binary_factorization,thm:peps_chrono_restriction}
    Let $U:\mathcal H_a\to\mathcal H_b$ be a source-free composition, and let
    $\mathbf p$ be a list of distinct parties containing every party occurring
    in $U$, including parties named by scalar local maps. For the canonical
    grouping isometries $J_a,J_b$ associated with this order,
    \begin{align}
        J_b U=\left(\bigotimes_{p\in\mathbf p}U|_p\right)J_a.
        \label{eq:peps_chrono_canonical_factors}
    \end{align}
    The factors are the actual restrictions of $U$. Zero local maps are permitted.
\end{theorem}

\begin{proof}
    \leanok
    Induct on the party list. Split off its first party using the binary
    factorization theorem, and apply induction to the complementary restriction.
    Theorem~\ref{thm:peps_chrono_restriction} identifies the remaining singleton
    factors with the original restrictions. With an empty party list there are
    no participating operations, so the composition is the scalar identity.
\end{proof}

\begin{theorem}[Complete chronological factorization]
    \label{thm:peps_chrono_complete_factors}
    \lean{TNLean.PEPS.PairEffect.SourceCircuit.groupByPartyIso_fullBranchWord_eq_tensorLocalRemainingMaps,
        TNLean.PEPS.PairEffect.SourceCircuit.groupByPartyIso_fullRemainingWord_eq_tensorLocalRemainingMaps}
    \leanok
    \uses{thm:peps_chrono_canonical_factors,def:peps_chrono_complete_local_maps,thm:peps_chrono_complete_preparation}
    Suppose $\mathbf p$ lists all parties exactly once. Then
    \begin{align}
        J_b V_\lambda
         & =\left(\bigotimes_{p\in\mathbf p}L_p(\pi_{\{p\}}\lambda)\right)
        J_{\mathcal S_W\mathbin{+}a},
        \label{eq:peps_chrono_complete_factors}                            \\
        J_b W_\lambda x
         & =\left(\bigotimes_{p\in\mathbf p}L_p(\pi_{\{p\}}\lambda)\right)
        J_{\mathcal S_W\mathbin{+}a}(s_\lambda\otimes x).
        \label{eq:peps_chrono_complete_vectors}
    \end{align}
    No allowedness assumption is needed for these algebraic identities.
\end{theorem}

\begin{proof}
    \leanok
    Apply (\ref{eq:peps_chrono_canonical_factors}) to the source-free residual,
    substitute the defining local maps, and then use
    (\ref{eq:peps_chrono_complete_preparation}).
\end{proof}

\subsection{The actual all-party partial expansion}

The all-party partial expansion differs from a full label assignment at
empty-participant gates: it retains each such gate as its original aggregate
operator. Denote its actual fixed-source residual by $R_\xi$, for
$\xi\in\mathcal C_P(W)$. Its owners are the original parties together with
one auxiliary exterior owner $\ast$. Every original register is owned by
its original party, and no original register is owned by $\ast$.
Let $\mathcal S$ be its fixed free-source list, so its input is
$\mathcal S\mathbin{+}a$ after this owner relabelling.

\begin{theorem}[Incident choices determine the actual partial residual]
    \label{thm:peps_chrono_partial_locality}
    \lean{TNLean.PEPS.PairEffect.SourceCircuit.restrict_partialResidual_eval_eq_of_projectAllChoices_eq,
        TNLean.PEPS.PairEffect.SourceCircuit.restrict_partialWord_localPart_eval_heq_of_projectAllChoices_eq}
    \leanok
    \uses{thm:peps_chrono_disjoint,thm:peps_chrono_restriction,def:peps_chrono_full_labels}
    Choose any subset $A^+$ of the enlarged owner set, and let $A$ be its inverse
    image in the original parties. If
    $\pi_A^P\xi=\pi_A^P\zeta$, then $R_\xi|_{A^+}=R_\zeta|_{A^+}$.
    The same equality holds before the final identification with the fixed
    source registers, for the actual source-free parts of the partial branches.
\end{theorem}

\begin{proof}
    \leanok
    Induct on the original chronology. At a touched gate the selected labels
    are equal. A gate disjoint from $A$ restricts to the identity on its selected
    spectators, even after moving sources to the beginning. An empty-participant
    gate has the same aggregate operator in both branches. The restriction
    identities propagate these equalities through compositions, spectators, and
    the final fixed-source identifications.
\end{proof}

\begin{definition}[Actual partial local residuals]
    \label{def:peps_chrono_partial_local_maps}
    \lean{TNLean.PEPS.PairEffect.SourceCircuit.allResidualInput,
        TNLean.PEPS.PairEffect.SourceCircuit.norm_partialLocalRemainingMap_le_one,
        TNLean.PEPS.PairEffect.SourceCircuit.partialLocalRemainingMap,
        TNLean.PEPS.PairEffect.SourceCircuit.restrict_partialResidual_eq_partialLocalRemainingMap}
    \leanok
    \uses{thm:peps_chrono_partial_locality,def:peps_canonical_source_residual}
    For $\eta\in\mathcal C_A(W)$ define $B_{A^+}(\eta)$ by restricting
    $R_\xi$ for an all-party partial choice extending $\eta$, and set it to zero
    if no extension exists. The common domain is the restriction of the fixed
    input $\mathcal S\mathbin{+}a$. Theorem~\ref{thm:peps_chrono_partial_locality}
    gives
    $R_\xi|_{A^+}=B_{A^+}(\pi_A^P\xi)$ for every actual $\xi$.
    If $W$ is allowed, then $\|B_{A^+}(\eta)\|\le1$ for every $\eta$.
\end{definition}

\begin{proof}
    \leanok
    Any two extending choices have identical restricted evaluations. Allowedness
    of the actual residual and of its restriction gives the norm bound; the
    unextendable case is the zero map.
\end{proof}

\begin{theorem}[Factorization of the actual partial residual]
    \label{thm:peps_chrono_partial_factors}
    \lean{TNLean.PEPS.PairEffect.SourceCircuit.groupByPartyIso_partialResidual_eq_tensorPartialLocalRemainingMaps}
    \leanok
    \uses{thm:peps_chrono_canonical_factors,def:peps_chrono_partial_local_maps}
    Let $\mathbf q$ list all enlarged owners exactly once, including $\ast$.
    For the induced singleton choices $\xi_q$,
    \begin{align}
        J_b R_\xi=\left(\bigotimes_{q\in\mathbf q}B_{\{q\}}(\xi_q)\right)
        J_{\mathcal S\mathbin{+}a}.
        \label{eq:peps_chrono_partial_factors}
    \end{align}
    The auxiliary scalar factor remains present. The identity also permits an
    empty original party set.
\end{theorem}

\begin{proof}
    \leanok
    Apply the canonical singleton-restriction factorization to $R_\xi$ and use
    the equality defining its actual partial local maps.
\end{proof}

\begin{theorem}[Evaluation with spectator registers]
    \label{thm:peps_chrono_spectators}
    \lean{TNLean.PEPS.PairEffect.Word.eval_frameList_tmul}
    \leanok
    Let $U:\mathcal H_a\to\mathcal H_b$ be a composition and let $r$ be any
    ordered spectator list. Retaining these registers gives
    $(I_r\otimes U)(x\otimes y)=x\otimes Uy$ for every
    $x\in\mathcal H_r$ and $y\in\mathcal H_a$, under the canonical
    concatenation isometries.
\end{theorem}

\begin{proof}
    \leanok
    Induct on the spectator list. The empty list is the scalar factor and the
    assertion is linearity. For a nonempty list, elementary tensors span its
    memory; tensoring the induction hypothesis with the first register proves
    the assertion.
\end{proof}

\begin{definition}[Preparing the original private input]
    \label{def:peps_chrono_private_input}
    \lean{TNLean.PEPS.PairEffect.SourceCircuit.allResidualSourceLayout,
        TNLean.PEPS.PairEffect.SourceCircuit.partialSourceWord,
        TNLean.PEPS.PairEffect.SourceCircuit.sourcePrivatePreparation,
        TNLean.PEPS.PairEffect.SourceCircuit.sources_partialSourceWord,
        TNLean.PEPS.PairEffect.SourceCircuit.sources_sourcePrivatePreparation,
        TNLean.PEPS.PairEffect.SourceCircuit.eval_partialSourceWord_tmul,
        TNLean.PEPS.PairEffect.SourceCircuit.eval_sourcePrivatePreparation_tmul}
    \leanok
    \uses{thm:peps_chrono_spectators,def:peps_chrono_partial_local_maps}
    Fix the original normalized private product input $\psi\in\mathcal H_a$.
    Its local preparation, retaining every source register, is a source-free
    composition $P_\psi:\mathcal H_{\mathcal S}\to
        \mathcal H_{\mathcal S\mathbin{+}a}$ satisfying
    $P_\psi y=y\otimes\psi$ for arbitrary $y$.
    Define $D_\xi=R_\xi P_\psi$. It is an actual source-free composition and
    \begin{align}
        D_\xi y=R_\xi(y\otimes\psi).
        \label{eq:peps_chrono_private_input}
    \end{align}
    The scalar terminal factor of the empty input is $1$. All owner and
    concatenation identifications in this formula are canonical.
\end{definition}

\begin{proof}
    \leanok
    Prepare the original private product vector by its recorded local maps and
    retain the source registers as spectators. Use the spectator evaluation
    identity and owner-relabelling naturality. Neither this preparation nor the
    actual residual introduces a pair source.
\end{proof}

\begin{definition}[The actual local source-to-output maps]
    \label{def:peps_chrono_source_local_maps}
    \lean{TNLean.PEPS.PairEffect.SourceCircuit.partialSourceLocalMap,
        TNLean.PEPS.PairEffect.SourceCircuit.restrict_partialSourceWord_eq_partialSourceLocalMap}
    \leanok
    \uses{def:peps_chrono_private_input,def:peps_chrono_partial_local_maps}
    For a selected owner set $A^+$ and its induced original party set $A$, put
    $D_{A^+}(\eta)=B_{A^+}(\eta)(P_\psi|_{A^+})$.
    Then $D_\xi|_{A^+}=D_{A^+}(\pi_A^P\xi)$.
    Thus the actual map from the source registers of these owners to their
    output depends only on the incident chronological choices.
\end{definition}

\begin{proof}
    \leanok
    Restrict the composition $R_\xi P_\psi$ and substitute the incident-choice
    identity for the first factor. The private preparation is fixed independently
    of the chronological labels.
\end{proof}

\begin{theorem}[The actual local maps factor the prepared chronology]
    \label{thm:peps_chrono_source_factors}
    \lean{TNLean.PEPS.PairEffect.SourceCircuit.groupByPartyIso_partialSourceWord_eq_tensorPartialSourceLocalMaps}
    \leanok
    \uses{thm:peps_chrono_canonical_factors,def:peps_chrono_source_local_maps}
    With the same exhaustive distinct order $\mathbf q$ of enlarged owners,
    \begin{align}
        J_b D_\xi=\left(\bigotimes_{q\in\mathbf q}D_{\{q\}}(\xi_q)\right)J_{\mathcal S}.
        \label{eq:peps_chrono_source_factors}
    \end{align}
    The factor for $\ast$ is included. This is an identity for every free source
    input, before replacing the source operators or tracing private outputs.
\end{theorem}

\begin{proof}
    \leanok
    Apply the canonical singleton-restriction factorization to the source-free
    composition $D_\xi$, and identify each restriction with its actual local map.
\end{proof}

\begin{definition}[Incident labels carried by the gate stars]
    \label{def:peps_chrono_incident_labels}
    \lean{TNLean.PEPS.PairEffect.SourceCircuit.OriginalGateLabelsAt,
        TNLean.PEPS.PairEffect.SourceCircuit.choicesOfOriginalGateLabelsAt,
        TNLean.PEPS.PairEffect.SourceCircuit.choicesOfOriginalGateLabelsAt_eq_projectAllChoices,
        TNLean.PEPS.PairEffect.SourceCircuit.isTouched_groupedParty_iff,
        TNLean.PEPS.PairEffect.SourceCircuit.originalGateLabelsAtOfAllChoices,
        TNLean.PEPS.PairEffect.SourceCircuit.partySelector_affectedOwner_all_apply}
    \leanok
    \uses{def:peps_chrono_touched_labels,def:peps_chrono_full_labels,def:compression_operator_gate_labels}
    At an original party $p$, the incident labels are precisely
    $(\xi_g)_{g\in G_+(W),\ p\in P_g}$.
    Selecting the enlarged owner corresponding to $p$ touches exactly these
    original occurrences. The incident labels therefore reconstruct a partial
    choice at $p$. For an actual all-party choice, restricting its labels to these
    stars and reconstructing the partial choice returns $\pi_{\{p\}}^P\xi$.
    No unused gate label is chosen.
\end{definition}

\begin{proof}
    \leanok
    The owner relabelling preserves the original party label, so its singleton
    selector touches a gate exactly when that party participates. Apply the
    inverse of (\ref{eq:peps_chrono_touched_labels}) and the label-recovery
    identity of the all-party projection.
\end{proof}

\begin{definition}[Local source matrices and private traces]
    \label{def:peps_chrono_local_matrices}
    \lean{TNLean.PEPS.PairEffect.SourceCircuit.partialSourceLocalMatrix,
        TNLean.PEPS.PairEffect.SourceCircuit.partialSourceLocalMatrix_eq_actual_restriction,
        TNLean.PEPS.PairEffect.SourceCircuit.partialSourceLocalPhysicalMatrix}
    \leanok
    \uses{def:peps_chrono_source_local_maps}
    Choose finite orthonormal bases $b_i$ of the selected source memory and
    $c_j$ of the selected output memory. The actual local source matrix is
    $K_\eta(j,i)=\langle c_j,D_{A^+}(\eta)b_i\rangle$.
    For an incident choice obtained from a genuine all-party choice, these are
    exactly the coordinates of $D_\xi|_{A^+}$.
    If the output basis is indexed by physical and private coordinates
    $(x,e)$, its mixed local physical matrix for any source matrix $X$ is
    \begin{align}
        M_{\eta,\zeta}(X)=\operatorname{Tr}_{\mathrm{priv}}
        (K_\eta X K_\zeta^\dagger).
        \label{eq:peps_chrono_local_matrices}
    \end{align}
    The source matrix and the two incident choices are arbitrary; positivity
    is not asserted.
\end{definition}

\begin{proof}
    \leanok
    Take matrix coordinates of the actual restricted map equality. The physical
    matrix is its stated partial trace, with the ket map on the left and the
    adjoint bra map on the right.
\end{proof}

\subsection{Empty-participant gates and their scalar factor}
\begin{theorem}[A gate with no participants]
    \label{thm:peps_chrono_empty_gate}
    \lean{TNLean.PEPS.PairEffect.PreparedSourceGate.eval_heq_coeffSum_smul_id_of_isEmpty,
        TNLean.PEPS.PairEffect.PreparedSourceGate.norm_coeffSum_le_one_of_isEmpty,
        TNLean.PEPS.PairEffect.SourceCircuit.eval_gate_eq_coeffSum_smul_id_of_isEmpty,
        TNLean.PEPS.PairEffect.SourceCircuit.eval_gate_heq_coeffSum_smul_id_of_participants_eq_empty}
    \leanok
    \uses{def:peps_prepared_source_gate}
    If a prepared gate has no participants, its actual aggregate map is
    $\sigma_g I$, where $\sigma_g=\sum_{\xi\in\Xi_g}c_g(\xi)$ and
    $|\sigma_g|\le1$. This holds beside arbitrary spectator registers.
    The same conclusion follows when the literal participant set of the placed
    gate is empty. An empty branch set gives the zero scalar.
\end{theorem}

\begin{proof}
    \leanok
    A source-free word on an empty party set contains no local map, source, or
    register exchange and evaluates to the scalar identity. Apply this to every
    prepared branch and sum its original coefficients. The allowed aggregate
    gate gives the modulus bound. Injectivity of its owner map identifies an
    empty literal participant set with an empty local party type.
\end{proof}

\begin{definition}[The product of empty-participant gate scalars]
    \label{def:peps_chrono_empty_scalar}
    \lean{TNLean.PEPS.PairEffect.SourceCircuit.emptyGateScalar,
        TNLean.PEPS.PairEffect.SourceCircuit.emptyGateScalar_gate_of_isEmpty,
        TNLean.PEPS.PairEffect.SourceCircuit.eval_heq_emptyGateScalar_smul_id_of_isEmpty,
        TNLean.PEPS.PairEffect.SourceCircuit.norm_emptyGateScalar_le_one}
    \leanok
    \uses{thm:peps_chrono_empty_gate}
    Define
    \begin{align}
        \sigma(W)=\prod_{g\in G(W),\ P_g=\varnothing}
        \left(\sum_{\xi\in\Xi_g}c_g(\xi)\right).
        \label{eq:peps_chrono_empty_scalar}
    \end{align}
    Equivalently, identity, private-map, and exchange operations contribute $1$;
    composition multiplies the scalars, and adjoining a spectator leaves them
    unchanged. A gate contributes its coefficient sum exactly when its participant
    type is empty. Every chronological composition satisfies $|\sigma(W)|\le1$.
    If the original party set is empty, its actual evaluation is $\sigma(W)I$.
\end{definition}

\begin{proof}
    \leanok
    Induct on the chronology. The only nontrivial factor is the single empty gate
    of Theorem~\ref{thm:peps_chrono_empty_gate}. Multiplication preserves the
    modulus bound. On an empty party set all remaining possible operations are
    identities and compositions of these scalar gates.
\end{proof}

\begin{theorem}[The actual exterior scalar in the all-party expansion]
    \label{thm:peps_chrono_none_scalar}
    \lean{TNLean.PEPS.PairEffect.SourceCircuit.restrict_partialWord_localPart_gate_heq_coeffSum_of_isEmpty,
        TNLean.PEPS.PairEffect.SourceCircuit.restrict_sources_partialWord_all_eq_nil,
        TNLean.PEPS.PairEffect.Word.restrict_localPart_comp_scalar,
        TNLean.PEPS.PairEffect.Word.restrict_localPart_frame_scalar,
        TNLean.PEPS.PairEffect.SourceCircuit.restrict_partialResidual_all_eval_heq_emptyGateScalar,
        TNLean.PEPS.PairEffect.SourceCircuit.restrict_partialWord_all_localPart_eval_heq_emptyGateScalar}
    \leanok
    \uses{def:peps_chrono_empty_scalar,thm:peps_chrono_disjoint,thm:peps_chrono_restriction}
    No source register of an all-party partial branch has the auxiliary owner
    $\ast$. Restricting either its source-free part or its fixed-source residual
    to this owner gives
    \begin{align}
        R_\xi|_{\{\ast\}}=\sigma(W)I_{\mathbb C}.
        \label{eq:peps_chrono_none_scalar}
    \end{align}
    At a single empty-participant gate the scalar is its original coefficient
    sum. Under composition the two scalars multiply. Retaining a register outside
    a selected owner set leaves its restricted scalar unchanged, provided the
    selected input and source register lists are empty.
\end{theorem}

\begin{proof}
    \leanok
    Every all-party relabelled register belongs to an original party, and every
    retained source has such owners. Induct on the actual chronology. Nonempty
    gates restrict to the identity at $\ast$; empty gates contribute their original
    coefficient sum. The restriction identities prove the composition and
    spectator steps, including the identifications with the fixed source list.
\end{proof}

\begin{theorem}[Absorbing a scalar at an existing party]
    \label{thm:peps_chrono_scalar_absorption}
    \lean{TNLean.PEPS.DependentBondNetwork.network_mul_at_vertex}
    \leanok
    Let a finite tensor network have a specified vertex $p$ and local tensors
    $T_q$. Replacing $T_p$ by $sT_p$ and keeping every other tensor and every
    bond space unchanged multiplies its contraction by $s$.
    There is no requirement that $s$ be nonzero.
\end{theorem}

\begin{proof}
    \leanok
    For each bond assignment, its local product contains the factor at $p$
    once. Extract $s$ from that finite product and then from the sum over bond
    assignments.
\end{proof}

\subsection{Operations on the actual source registers}
\begin{definition}[Separate operations on each source half]
    \label{def:peps_chrono_endpoint_word}
    \lean{TNLean.PEPS.PairEffect.SourceInventory.endpointWord,
        TNLean.PEPS.PairEffect.SourceInventory.sources_endpointWord,
        TNLean.PEPS.PairEffect.Word.localEndpoint,
        TNLean.PEPS.PairEffect.Word.sources_localEndpoint,
        TNLean.PEPS.PairEffect.Word.eval_localEndpoint_tmul}
    \leanok
    Let $R$ be an ordered pair-source list with halfspaces $U_i,V_i$ and actual
    owners $p_i,q_i$. For arbitrary continuous linear maps
    $A_i:U_i\to U_i$ and $B_i:V_i\to V_i$, form the actual source-free word that
    applies $A_i$ at $p_i$ and $B_i$ at $q_i$, in the original register order.
    A single register operation acts by $x\otimes y\mapsto A_i x\otimes y$,
    retaining the other registers as spectators. These operations introduce no
    pair source.
\end{definition}

\begin{proof}
    \leanok
    Use the local map on a singleton register and the identity on its remaining
    memory. Recursing over the pair list applies the two halfspace maps and then
    the remaining pairs, preserving their owners and order.
\end{proof}

\begin{theorem}[Local dependence of the source-half operations]
    \label{thm:peps_chrono_endpoint_locality}
    \lean{TNLean.PEPS.PairEffect.SourceInventory.restrict_endpointWord_eval_eq_of_agree,
        TNLean.PEPS.PairEffect.Word.restrict_localEndpoint_eval_eq_of_agree}
    \leanok
    \uses{def:peps_chrono_endpoint_word,thm:peps_chrono_restriction}
    Fix selected parties $A$. If $A_i=A_i'$ whenever $p_i\in A$, and
    $B_i=B_i'$ whenever $q_i\in A$, then the two source-half words have equal
    restricted evaluations on $A$. No equality is required for operators on
    unselected halves.
\end{theorem}

\begin{proof}
    \leanok
    For a singleton register, its operation is retained precisely when its owner
    is selected; otherwise its restriction is the identity. Induct on the pair
    list and use the composition and spectator restriction identities. The
    canonical register-list identifications preserve these equalities.
\end{proof}

\subsection{Coordinates of the party tensor product}
\label{sec:party-tensor-coordinates}

The coordinate identities below describe the local tensors and the local
summation over discarded registers in
\cite[proof of Theorem~5.2]{OpenAI2026PolynomialPEPS}.

\begin{definition}[Ordered tensor bases]
    \label{def:party-list-basis}
    \lean{TNLean.PEPS.PairEffect.partyListBasis,
        TNLean.PEPS.PairEffect.partyListBasis_nil_apply,
        TNLean.PEPS.PairEffect.partyListBasis_repr_nil,
        TNLean.PEPS.PairEffect.partyListBasis_cons_apply,
        TNLean.PEPS.PairEffect.partyListBasis_repr_cons_tmul,
        TNLean.PEPS.PairEffect.partyListBasis_repr_tensorPartyMaps}
    \leanok
    For every party $p$, let $H_p$ and $K_p$ be complex Hilbert spaces with
    finite orthonormal bases $(b_{p,i})_{i\in I_p}$ and
    $(c_{p,j})_{j\in J_p}$. Fix a finite ordered list
    $\mathbf p=(p_0,\ldots,p_{m-1})$, allowing repeated parties. For $\mathbf i\in\prod_{0\le r<m}I_{p_r}$, its tensor basis is
    \begin{align}
        b_{\mathbf p,\mathbf i}
        =b_{p_0,i_0}\otimes\cdots\otimes b_{p_{m-1},i_{m-1}},
    \end{align}
    with the prescribed order and association of the tensor factors. The empty
    product is $\mathbb C$ with its basis vector $1$; the coordinate of a scalar
    $z$ in this basis is $z$. Adjoining a first party adjoins its basis vector as
    the first tensor factor. The coordinate of a simple tensor $u\otimes v$
    is the product of its two coordinates.

    For arbitrary continuous complex-linear maps $B_p:H_p\to K_p$, the matrix
    of their ordered tensor product satisfies
    \begin{align}
        \left[\left(\bigotimes_{0\le r<m} B_{p_r}\right)
            b_{\mathbf p,\mathbf i}\right]_{\mathbf j}
        =\prod_{0\le r<m} [B_{p_r}b_{p_r,i_r}]_{j_r}.
    \end{align}
    There is one factor for each occurrence in the list; equal party labels do
    not identify distinct positions.
\end{definition}

\begin{proof}
    \leanok
    Construct the basis by adjoining the first orthonormal basis to the tensor
    basis of the remaining list. The empty basis has the stated scalar
    coordinates. The coordinate identity for a simple tensor and the action of
    a tensor product of maps then give the product formula by induction on the
    length of the list.
\end{proof}

\begin{definition}[Coordinates indexed by the original parties]
    \label{def:party-labelled-tensor-coordinates}
    \lean{TNLean.PEPS.PairEffect.partyLabelledBasis,
        TNLean.PEPS.PairEffect.partyLabelledBasis_apply,
        TNLean.PEPS.PairEffect.partyLabelledBasis_repr_tensorPartyMaps}
    \uses{def:party-list-basis}
    \leanok
    Suppose the party set $P$ is finite and the ordered list $\mathbf p$
    contains every party exactly once. The preceding basis is then indexed by
    assignments $i\in\prod_{p\in P}I_p$: its vector at $i$ is the ordered tensor
    of the vectors $b_{p,i(p)}$. In the analogous output basis, for every
    $j\in\prod_{p\in P}J_p$,
    \begin{align}
        \left[\left(\bigotimes_{p\in\mathbf p}B_p\right)b_i\right]_j
        =\prod_{p\in P}[B_pb_{p,i(p)}]_{j(p)}.
    \end{align}
    The assertion includes the empty party set.
\end{definition}

\begin{proof}
    \leanok
    The position-to-party map of the exhaustive list is a bijection. Use it to
    relabel both tensor bases, then change the index in the finite product of
    local matrix entries.
\end{proof}

\begin{theorem}[Local summation of discarded coordinates]
    \label{thm:party-product-partial-trace}
    \lean{TNLean.PEPS.PairEffect.partialTraceRight_partyProductMatrix}
    \leanok
    Let $P$ be finite, let $X_p$ be sets, and let $E_p$ be finite sets. For
    arbitrary complex matrices $M_p$ indexed by $X_p\times E_p$, define the
    matrix $M$ indexed by
    $(\prod_p X_p)\times(\prod_p E_p)$ by
    \begin{align}
        M((x,e),(y,f))=\prod_{p\in P}M_p((x_p,e_p),(y_p,f_p)).
    \end{align}
    Tracing all discarded coordinates gives the product of the local partial
    traces:
    \begin{align}
        (\operatorname{Tr}_{\prod_p E_p}M)(x,y)
        =\prod_{p\in P}\left(\sum_{e_p\in E_p}
        M_p((x_p,e_p),(y_p,e_p))\right).
    \end{align}
    No positivity or trace normalization is required. Each discarded coordinate
    is summed within its own party.
\end{theorem}

\begin{proof}
    \leanok
    Expand the partial trace as a sum over $e\in\prod_pE_p$ and distribute the
    finite product of sums. The empty-set cases follow from the same identity.
\end{proof}
